\chapter{Trigonometric Functions}

Trigonometry plays a small part in economics and a large part in the
prerequisites, for two reasons. The graphs of $\sin$, $\cos$ and $\tan$
belong to the catalogue of shapes one is expected to recognise, and they are
the standard test functions of calculus: many derivative and limit
exercises use them precisely because they oscillate and so expose careless
reasoning. Seasonality and business cycles supply the
occasional genuine application.

The chapter therefore covers what is actually needed and stops. The unit
circle and the handful of exact values; the three graphs with their periods,
amplitudes and asymptotes; the transformation form
$a\sin\bigl(b(x-c)\bigr)+d$, which is the machinery of \chref{functions}{3} applied to a
periodic function; the four identities that carry all the simplification
work; and how to write down \emph{every} solution of a trigonometric
equation rather than just the one a calculator returns.

\section{Core Rules}

\begin{rulebox}{Radians, and why}
An angle in \textbf{radians} is the arc length it cuts on the unit circle,
so a full turn is $2\pi$:
\[
  180^{\circ}=\pi,\qquad
  x_{\text{rad}}=x_{\deg}\cdot\frac{\pi}{180},\qquad
  x_{\deg}=x_{\text{rad}}\cdot\frac{180}{\pi} .
\]
Calculus uses radians exclusively, and for one reason: only in radians is
$\dfrac{\mathrm{d}}{\mathrm{d}x}\sin x=\cos x$. In degrees an unwanted
factor $\pi/180$ appears in every derivative.
\end{rulebox}

\begin{rulebox}{The unit circle definition}
For the point on the unit circle reached by turning through the angle $x$
from the positive $x$-axis, $\cos x$ is its first coordinate and $\sin x$
its second. Hence, immediately:
\[
  \sin^{2}x+\cos^{2}x=1,\qquad
  -1\le\sin x\le1,\qquad -1\le\cos x\le1,
\]
the first being Pythagoras. Also $\tan x=\dfrac{\sin x}{\cos x}$, the slope
of the ray, undefined where $\cos x=0$.

Signs by quadrant --- $\sin$ is the height, $\cos$ the width:
\[
  \begin{array}{c|cccc}
     & \text{I} & \text{II} & \text{III} & \text{IV}\\\hline
    \sin & + & + & - & -\\
    \cos & + & - & - & +\\
    \tan & + & - & + & -
  \end{array}
\]
\end{rulebox}

\begin{rulebox}{The values to know by heart}
\[
  \begin{array}{c|ccccc}
    x & 0 & \pi/6 & \pi/4 & \pi/3 & \pi/2\\\hline
    \sin x & 0 & \tfrac12 & \tfrac{\sqrt2}{2} & \tfrac{\sqrt3}{2} & 1\\[0.3ex]
    \cos x & 1 & \tfrac{\sqrt3}{2} & \tfrac{\sqrt2}{2} & \tfrac12 & 0\\[0.3ex]
    \tan x & 0 & \tfrac{1}{\sqrt3} & 1 & \sqrt3 & \text{---}
  \end{array}
\]
Everything else follows from three facts: $\sin$ is \textbf{odd}
($\sin(-x)=-\sin x$), $\cos$ is \textbf{even} ($\cos(-x)=\cos x$), and both
have period $2\pi$ while $\tan$ has period $\pi$. So
$\sin\tfrac{5\pi}{6}=\sin\tfrac{\pi}{6}=\tfrac12$ (second quadrant, $\sin$
positive) and $\cos\tfrac{4\pi}{3}=-\cos\tfrac{\pi}{3}=-\tfrac12$ (third
quadrant, $\cos$ negative).
\end{rulebox}

\begin{gallery}
\begin{tikzpicture}
  \begin{axis}[fnaxis, xmin=-0.6, xmax=7.0, ymin=-1.6, ymax=1.6,
      xtick={3.1416, 6.2832}, xticklabels={$\pi$, $2\pi$},
      ytick={-1,1}, title={$y=\sin x$}]
    \addplot[domain=-0.4:6.9, samples=201]{sin(deg(x))};
  \end{axis}
\end{tikzpicture}\hfill
\begin{tikzpicture}
  \begin{axis}[fnaxis, xmin=-0.6, xmax=7.0, ymin=-1.6, ymax=1.6,
      xtick={3.1416, 6.2832}, xticklabels={$\pi$, $2\pi$},
      ytick={-1,1}, title={$y=\cos x$}]
    \addplot[domain=-0.4:6.9, samples=201]{cos(deg(x))};
  \end{axis}
\end{tikzpicture}\hfill
\begin{tikzpicture}
  \begin{axis}[fnaxis, xmin=-3.4, xmax=3.4, ymin=-3.4, ymax=3.4,
      xtick={-1.5708, 1.5708}, xticklabels={$-\frac{\pi}{2}$,$\frac{\pi}{2}$},
      ytick={-2,2}, title={$y=\tan x$}]
    \addplot[domain=-1.32:1.32, samples=101]{tan(deg(x))};
    \addplot[domain=1.82:3.3, samples=101]{tan(deg(x))};
    \addplot[domain=-3.3:-1.82, samples=101]{tan(deg(x))};
    \addplot[gray!55, dashed, thin, samples=2, forget plot]
      coordinates {(1.5708,-3.3)(1.5708,3.3)};
    \addplot[gray!55, dashed, thin, samples=2, forget plot]
      coordinates {(-1.5708,-3.3)(-1.5708,3.3)};
  \end{axis}
\end{tikzpicture}
\gcap{$\sin$ and $\cos$ are the same wave shifted by $\pi/2$, with period
$2\pi$ and range $[-1,1]$. $\tan$ has period $\pi$, range $\R$, and a
vertical asymptote at every odd multiple of $\pi/2$.}
\end{gallery}

\begin{rulebox}{Reading $a\sin\bigl(b(x-c)\bigr)+d$}
The four constants are exactly the four moves of \chref{functions}{3}, and each has a
name here:
\begin{itemize}[leftmargin=1.6em,itemsep=0.3ex]
  \item $|a|$ is the \textbf{amplitude}: the wave runs from $d-|a|$ to
        $d+|a|$, so the range is $[d-|a|,\,d+|a|]$. Negative $a$ flips it.
  \item $b$ sets the \textbf{period}, $T=\dfrac{2\pi}{|b|}$ --- larger $b$
        means faster oscillation. (For $\tan$, $T=\pi/|b|$.)
  \item $c$ is the \textbf{phase shift}: the whole wave moves $c$ to the
        right. Factor before reading: $\sin(2x-\pi)=\sin\bigl(2(x-\pi/2)\bigr)$
        is a shift of $\pi/2$, not of $\pi$.
  \item $d$ is the \textbf{midline}, the horizontal line the wave
        oscillates about.
\end{itemize}
So $y=3\sin\bigl(2(x-\tfrac{\pi}{4})\bigr)+1$ has amplitude $3$, period
$\pi$, midline $y=1$, range $[-2,4]$, shifted $\pi/4$ right.
\end{rulebox}

\begin{rulebox}{The identities that do the work}
\[
  \sin^{2}x+\cos^{2}x=1,\qquad
  \tan x=\frac{\sin x}{\cos x},\qquad
  1+\tan^{2}x=\frac{1}{\cos^{2}x} .
\]
\[
  \sin(x\pm y)=\sin x\cos y\pm\cos x\sin y,\qquad
  \cos(x\pm y)=\cos x\cos y\mp\sin x\sin y .
\]
Setting $y=x$ in the addition formulas gives the \textbf{double-angle}
forms, which are the ones that appear in practice:
\[
  \sin 2x=2\sin x\cos x,\qquad
  \cos 2x=\cos^{2}x-\sin^{2}x=1-2\sin^{2}x=2\cos^{2}x-1 .
\]
Almost every simplification is one of: replace $1-\cos^{2}x$ by
$\sin^{2}x$, split a $\tan$ into $\sin/\cos$, or collapse a double angle.
\end{rulebox}

\begin{rulebox}{Solving trigonometric equations: give \emph{all} solutions}
A calculator returns one angle; periodicity supplies infinitely many.
\begin{itemize}[leftmargin=1.6em,itemsep=0.35ex]
  \item $\sin x=s$ with $|s|\le1$: two families,
        $x=x_{0}+2k\pi$ and $x=\pi-x_{0}+2k\pi$, where
        $x_{0}=\arcsin s$ and $k\in\mathbb{Z}$.
  \item $\cos x=s$: $x=\pm x_{0}+2k\pi$ with $x_{0}=\arccos s$.
  \item $\tan x=s$: one family, $x=\arctan s+k\pi$, since the period is
        $\pi$.
  \item If the argument is $bx$ rather than $x$, solve for $bx$ first and
        divide \emph{everything}, the $2k\pi$ included: $\sin 2x=1$ gives
        $2x=\tfrac{\pi}{2}+2k\pi$, hence
        $x=\tfrac{\pi}{4}+k\pi$.
\end{itemize}
\end{rulebox}

\begin{rulebox}{Inverse trigonometric functions}
None of the three is injective, so each is inverted only on a restricted
domain:
\[
  \arcsin:[-1,1]\to\Bigl[-\tfrac{\pi}{2},\tfrac{\pi}{2}\Bigr],\qquad
  \arccos:[-1,1]\to[0,\pi],\qquad
  \arctan:\R\to\Bigl(-\tfrac{\pi}{2},\tfrac{\pi}{2}\Bigr).
\]
So $\arcsin(\sin x)=x$ only for $x$ in that range: $\arcsin(\sin\pi)=0$,
not $\pi$. And $\arctan$ is bounded, with horizontal asymptotes
$y=\pm\pi/2$ --- the reason it appears as a saturating S-shaped curve.
\end{rulebox}

\section{Worked Example}

\begin{worked}{Solve $2\sin^{2}x-\sin x=0$ on $[0,2\pi]$ and in general, and read off the shape of $y=3\sin\!\bigl(2(x-\tfrac{\pi}{4})\bigr)+1$}
\wstep{Step 1. Factor. Do not divide.}
\[
  2\sin^{2}x-\sin x=0
  \quad\Longrightarrow\quad
  \sin x\,(2\sin x-1)=0 .
\]
\wtrap{Dividing both sides by $\sin x$ gives $2\sin x-1=0$ and loses every
solution with $\sin x=0$ --- three of the five on $[0,2\pi]$. You may divide
by a quantity only when you know it is non-zero, and $\sin x$ is exactly the
quantity in question. Factor and use the zero-product property instead.}

\wstep{Step 2. Solve each factor on the given interval.}
\emph{First factor:} $\sin x=0$ on $[0,2\pi]$ gives $x=0,\ \pi,\ 2\pi$.

\emph{Second factor:} $\sin x=\tfrac12$. The reference angle is
$\pi/6$, and the sine is positive in the first and second quadrants, so
\[
  x=\frac{\pi}{6}\qquad\text{and}\qquad x=\pi-\frac{\pi}{6}=\frac{5\pi}{6}.
\]
\wtrap{$\arcsin(\tfrac12)=\pi/6$ is one solution, not the solution. The
inverse sine returns a single value from $[-\pi/2,\pi/2]$; the second
solution in each period must be supplied by hand from the symmetry
$\sin(\pi-x)=\sin x$. For cosine the partner is $-x$ (or $2\pi-x$), for
tangent it is $x+\pi$.}

\wstep{Step 3. Collect the five solutions.}
\[
  \boxed{\;x\in\Bigl\{0,\ \tfrac{\pi}{6},\ \tfrac{5\pi}{6},\ \pi,\ 2\pi\Bigr\}\;}
\]
Check the smallest and one middle value: at $x=0$, $0-0=0$ \checkmark; at
$x=\pi/6$, $2\cdot\tfrac14-\tfrac12=0$ \checkmark.

\wstep{Step 4. Give the general solution, since ``all solutions'' usually means all of $\R$.}
\[
  x=k\pi,
  \qquad
  x=\frac{\pi}{6}+2k\pi,
  \qquad
  x=\frac{5\pi}{6}+2k\pi,
  \qquad k\in\mathbb{Z} .
\]
Note the two different periods: the zeros of $\sin$ repeat every $\pi$
(because $\sin(x+\pi)=-\sin x$, which is also zero), while the solutions of
$\sin x=\tfrac12$ repeat only every $2\pi$.
\wtrap{Adding $2k\pi$ to everything is safe but incomplete for the first
family; writing $k\pi$ for the second family is simply wrong. Work out the
period of each family separately from the equation it solves.}

\wstep{Step 5. Now the sinusoid: read the four parameters off $a\sin\bigl(b(x-c)\bigr)+d$.}
For $y=3\sin\bigl(2(x-\tfrac{\pi}{4})\bigr)+1$:
\begin{center}\footnotesize
\begin{tabular}{llll}
  \hline
  parameter & value & name & effect\\
  \hline
  $a$ & $3$ & amplitude & oscillates $3$ above and below the midline\\
  $b$ & $2$ & angular frequency & period $=2\pi/b=\pi$\\
  $c$ & $\pi/4$ & phase shift & shifted \emph{right} by $\pi/4$\\
  $d$ & $1$ & midline & centred on $y=1$, so range $[-2,4]$\\
  \hline
\end{tabular}
\end{center}

\wstep{Step 6. Sketch it from those four numbers.}
\begin{center}
\begin{tikzpicture}
\begin{axis}[fnaxis, width=8.2cm, height=4.4cm,
             domain=-0.4:3.6, samples=201,
             xmin=-0.5, xmax=3.7, ymin=-2.8, ymax=4.8,
             xtick={0.7854, 1.5708, 2.3562, 3.1416},
             xticklabels={$\frac{\pi}{4}$,$\frac{\pi}{2}$,$\frac{3\pi}{4}$,$\pi$},
             ytick={-2,1,4}]
  \addplot[accent!35, thick, dashed, samples=2, forget plot] {1};
  \addplot {3*sin(deg(2*(x-pi/4)))+1};
\end{axis}
\end{tikzpicture}
\end{center}
The curve crosses its midline $y=1$ going upward at $x=\pi/4$ --- that is
what the phase shift means --- peaks at $4$ a quarter-period later
($x=\pi/4+\pi/4=\pi/2$), and repeats after one full period $\pi$.
\wtrap{The period is $2\pi/b$, not $2\pi b$: a \emph{larger} $b$ gives a
\emph{shorter} period. And the phase shift is only readable once the inside
is factored as $b(x-c)$; from $\sin(2x-\tfrac{\pi}{2})$ the shift is
$\pi/4$, not $\pi/2$. It is the same trap as in \chref{functions}{3} on
transformations.}
\end{worked}

\section{Practice Problems}

Convert in problems 1--6.

\begin{problems}
  \item $60^{\circ}$ into radians.
  \item $270^{\circ}$ into radians.
  \item $\tfrac{3\pi}{4}$ into degrees.
  \item $\tfrac{5\pi}{3}$ into degrees.
  \item $1$ radian into degrees, to one decimal.
  \item Why does calculus insist on radians?
\end{problems}

Give the exact value in problems 7--16.

\begin{problems}[resume]
  \item $\sin\tfrac{\pi}{6}$
  \item $\cos\tfrac{\pi}{4}$
  \item $\tan\tfrac{\pi}{3}$
  \item $\sin\tfrac{5\pi}{6}$
  \item $\cos\tfrac{2\pi}{3}$
  \item $\tan\pi$
  \item $\sin\bigl(-\tfrac{\pi}{3}\bigr)$
  \item $\cos\tfrac{7\pi}{2}$
  \item $\sin\tfrac{9\pi}{4}$
  \item $\tan\tfrac{3\pi}{4}$
\end{problems}

In problems 17--24 give period, amplitude (where defined), range and any
asymptotes.

\begin{problems}[resume]
  \item $y=\sin x$
  \item $y=\cos x$
  \item $y=\tan x$
  \item $y=3\sin x$
  \item $y=\sin 2x$
  \item $y=\cos\tfrac{x}{2}$
  \item $y=2\sin x+1$
  \item $y=\tan 3x$
\end{problems}

\begin{problems}[resume]
  \item Describe $y=\sin\bigl(x-\tfrac{\pi}{2}\bigr)$ as a shift, and name
        the simpler function it equals.
  \item Give amplitude, period, midline and range of
        $y=3\sin\bigl(2(x-\tfrac{\pi}{4})\bigr)+1$.
  \item Give amplitude, period and phase shift of
        $y=\sin(2x-\pi)$ --- factor first.
  \item Give amplitude, period and range of $y=-2\cos 3x$.
  \item Write a sine function with amplitude $4$, period $\pi$ and midline
        $y=-1$.
  \item Write a cosine function with amplitude $2$, period $6\pi$ and no
        vertical shift.
  \item Monthly sales oscillate about $500$ units with a swing of $\pm80$
        and a twelve-month cycle, peaking in month $3$. Write $S(t)$.
  \item Why can no function of the form $a\sin(bx)+d$ ever be injective on
        $\R$?
\end{problems}

Simplify or verify in problems 33--38.

\begin{problems}[resume]
  \item $\dfrac{1-\cos^{2}x}{\sin x}$
  \item $\tan x\cdot\cos x$
  \item $\sin^{2}x-1$
  \item $\dfrac{\sin 2x}{2\cos x}$
  \item Verify $1+\tan^{2}x=\dfrac{1}{\cos^{2}x}$.
  \item Verify $\cos 2x=1-2\sin^{2}x$ from the addition formula.
\end{problems}

Give \emph{all} real solutions in problems 39--45.

\begin{problems}[resume]
  \item $\sin x=\tfrac12$
  \item $\cos x=0$
  \item $\tan x=1$
  \item $2\cos x+1=0$
  \item $\sin 2x=1$
  \item $\cos^{2}x=\tfrac14$
  \item $\sin x=\cos x$
\end{problems}

\section{Solutions}

\begin{solutions}
  \item $60\cdot\tfrac{\pi}{180}=\tfrac{\pi}{3}$.
  \item $270\cdot\tfrac{\pi}{180}=\tfrac{3\pi}{2}$.
  \item $\tfrac{3\pi}{4}\cdot\tfrac{180}{\pi}=135^{\circ}$.
  \item $\tfrac{5\pi}{3}\cdot\tfrac{180}{\pi}=300^{\circ}$.
  \item $\tfrac{180}{\pi}\approx57.3^{\circ}$.
  \item Because $\tfrac{\mathrm{d}}{\mathrm{d}x}\sin x=\cos x$ holds only
        in radians. In degrees the limit
        $\lim_{h\to0}\tfrac{\sin h}{h}$ equals $\tfrac{\pi}{180}$ instead
        of $1$, and that factor would appear in every derivative and integral
        of a trigonometric function.
  \item $\tfrac12$.
  \item $\tfrac{\sqrt2}{2}$.
  \item $\sqrt3$.
  \item Second quadrant, $\sin$ positive:
        $\sin\tfrac{5\pi}{6}=\sin\tfrac{\pi}{6}=\tfrac12$.
  \item Second quadrant, $\cos$ negative:
        $-\cos\tfrac{\pi}{3}=-\tfrac12$.
  \item $\tan\pi=\tfrac{\sin\pi}{\cos\pi}=\tfrac{0}{-1}=0$.
  \item $\sin$ is odd: $-\sin\tfrac{\pi}{3}=-\tfrac{\sqrt3}{2}$.
  \item $\tfrac{7\pi}{2}=\tfrac{3\pi}{2}+2\pi$, so
        $\cos\tfrac{7\pi}{2}=\cos\tfrac{3\pi}{2}=0$.
  \item $\tfrac{9\pi}{4}=\tfrac{\pi}{4}+2\pi$, so the value is
        $\sin\tfrac{\pi}{4}=\tfrac{\sqrt2}{2}$.
  \item Second quadrant, $\tan$ negative:
        $-\tan\tfrac{\pi}{4}=-1$.
  \item Period $2\pi$, amplitude $1$, range $[-1,1]$, no asymptotes.
  \item Period $2\pi$, amplitude $1$, range $[-1,1]$, no asymptotes.
  \item Period $\pi$, amplitude undefined (unbounded), range $\R$,
        vertical asymptotes at $x=\tfrac{\pi}{2}+k\pi$.
  \item Period $2\pi$, amplitude $3$, range $[-3,3]$.
  \item Period $\tfrac{2\pi}{2}=\pi$, amplitude $1$, range $[-1,1]$.
  \item Period $\tfrac{2\pi}{1/2}=4\pi$, amplitude $1$, range $[-1,1]$.
  \item Period $2\pi$, amplitude $2$, midline $y=1$, range $[-1,3]$.
  \item Period $\tfrac{\pi}{3}$, unbounded range $\R$, asymptotes where
        $3x=\tfrac{\pi}{2}+k\pi$, i.e. $x=\tfrac{\pi}{6}+\tfrac{k\pi}{3}$.
  \item A shift of $\tfrac{\pi}{2}$ to the right. It equals $-\cos x$,
        since $\sin\bigl(x-\tfrac{\pi}{2}\bigr)
        =\sin x\cos\tfrac{\pi}{2}-\cos x\sin\tfrac{\pi}{2}=-\cos x$.
  \item Amplitude $3$, period $\tfrac{2\pi}{2}=\pi$, midline $y=1$, range
        $[-2,4]$, shifted $\tfrac{\pi}{4}$ to the right.
  \item $\sin(2x-\pi)=\sin\bigl(2(x-\tfrac{\pi}{2})\bigr)$: amplitude $1$,
        period $\pi$, phase shift $\tfrac{\pi}{2}$ right. Reading the
        shift as $\pi$ from the unfactored form is the standard error.
  \item Amplitude $2$ (the minus sign reflects rather than changing the
        amplitude), period $\tfrac{2\pi}{3}$, range $[-2,2]$.
  \item Period $\pi$ needs $b=2$: $y=4\sin(2x)-1$.
  \item Period $6\pi$ needs $b=\tfrac13$:
        $y=2\cos\bigl(\tfrac{x}{3}\bigr)$.
  \item Midline $500$, amplitude $80$, period $12$ so
        $b=\tfrac{2\pi}{12}=\tfrac{\pi}{6}$, and a sine peaks a quarter
        period after its shift, so peaking at $t=3$ means $c=0$:
        \[
          S(t)=500+80\sin\!\left(\frac{\pi t}{6}\right).
        \]
        (Check: at $t=3$ the sine is $\sin\tfrac{\pi}{2}=1$, giving the
        peak $580$.)
  \item Because it is periodic: $f(x+T)=f(x)$ for $T=2\pi/|b|\neq0$, so two
        different inputs share a value. Injectivity fails everywhere, which
        is why the inverse functions need restricted domains.
  \item $\dfrac{\sin^{2}x}{\sin x}=\sin x$ (for $\sin x\neq0$).
  \item $\dfrac{\sin x}{\cos x}\cdot\cos x=\sin x$ (for
        $\cos x\neq0$).
  \item $-(1-\sin^{2}x)=-\cos^{2}x$.
  \item $\dfrac{2\sin x\cos x}{2\cos x}=\sin x$ (for $\cos x\neq0$).
  \item Divide $\sin^{2}x+\cos^{2}x=1$ by $\cos^{2}x$:
        $\tan^{2}x+1=\dfrac{1}{\cos^{2}x}$.
  \item $\cos2x=\cos^{2}x-\sin^{2}x$, and substituting
        $\cos^{2}x=1-\sin^{2}x$ gives $1-\sin^{2}x-\sin^{2}x
        =1-2\sin^{2}x$.
  \item $x=\tfrac{\pi}{6}+2k\pi$ or $x=\tfrac{5\pi}{6}+2k\pi$,
        $k\in\mathbb{Z}$.
  \item $x=\tfrac{\pi}{2}+k\pi$.
  \item $x=\tfrac{\pi}{4}+k\pi$.
  \item $\cos x=-\tfrac12$, so
        $x=\pm\tfrac{2\pi}{3}+2k\pi$.
  \item $2x=\tfrac{\pi}{2}+2k\pi$, hence
        $x=\tfrac{\pi}{4}+k\pi$.
  \item $\cos x=\pm\tfrac12$, giving
        $x=\pm\tfrac{\pi}{3}+2k\pi$ and $x=\pm\tfrac{2\pi}{3}+2k\pi$ ---
        more compactly $x=\tfrac{\pi}{3}+\tfrac{k\pi}{1}$ fails, so keep
        the four families, or write
        $x=\pm\tfrac{\pi}{3}+k\pi$.
  \item Dividing by $\cos x$ (legitimate, since $\cos x=0$ would force
        $\sin x=0$ too, impossible): $\tan x=1$, so
        $x=\tfrac{\pi}{4}+k\pi$.
\end{solutions}

\section{Additional Problems}

\begin{problems}
  \item $45^{\circ}$ and $150^{\circ}$ into radians.
  \item $\tfrac{7\pi}{6}$ and $\tfrac{11\pi}{6}$ into degrees.
  \item $\cos\tfrac{5\pi}{4}$
  \item $\sin\tfrac{4\pi}{3}$
  \item $\tan\tfrac{5\pi}{6}$
  \item $\sin\tfrac{11\pi}{6}$
  \item $\cos(-\tfrac{\pi}{6})$
  \item $\tan\tfrac{7\pi}{4}$
  \item Period, amplitude and range of $y=\tfrac12\cos 4x-2$.
  \item Period and asymptotes of $y=\tan\bigl(x-\tfrac{\pi}{4}\bigr)$.
  \item Amplitude, period, midline of
        $y=5\cos\bigl(\tfrac{\pi}{3}(x-2)\bigr)+3$.
  \item Write a cosine with maximum $7$, minimum $1$ and period $4$.
  \item Simplify $\dfrac{\cos^{2}x-1}{\cos x-1}$.
  \item Simplify $\sin x\cos x\cdot\dfrac{2}{\sin 2x}$.
  \item Simplify $\dfrac{1}{\cos^{2}x}-\tan^{2}x$.
  \item Verify $\sin(x+\pi)=-\sin x$.
  \item Verify $\cos\bigl(\tfrac{\pi}{2}-x\bigr)=\sin x$.
  \item Solve $\sin x=-\tfrac{\sqrt3}{2}$.
  \item Solve $\tan 2x=\sqrt3$.
  \item Solve $2\sin^{2}x-\sin x=0$.
  \item Solve $\cos 2x=\tfrac12$.
  \item Solve $\sin x=2$ and explain the answer.
  \item Give $\arcsin\tfrac12$, $\arccos(-1)$ and $\arctan(-1)$.
  \item Evaluate $\arcsin\bigl(\sin\tfrac{3\pi}{4}\bigr)$ and explain why
        the answer is not $\tfrac{3\pi}{4}$.
  \item Quarterly demand is $D(t)=200-40\cos\bigl(\tfrac{\pi t}{2}\bigr)$
        with $t$ in quarters. Give the period, the maximum and minimum
        demand, and the quarters in which each occurs.
\end{problems}

\section{Additional Solutions}

\begin{solutions}
  \item $\tfrac{\pi}{4}$ and $\tfrac{5\pi}{6}$.
  \item $210^{\circ}$ and $330^{\circ}$.
  \item Third quadrant, $\cos$ negative:
        $-\tfrac{\sqrt2}{2}$.
  \item Third quadrant, $\sin$ negative:
        $-\tfrac{\sqrt3}{2}$.
  \item Second quadrant, $\tan$ negative:
        $-\tfrac{1}{\sqrt3}$.
  \item Fourth quadrant, $\sin$ negative:
        $-\tfrac12$.
  \item $\cos$ is even: $\cos\tfrac{\pi}{6}=\tfrac{\sqrt3}{2}$.
  \item Fourth quadrant, $\tan$ negative: $-1$.
  \item Period $\tfrac{2\pi}{4}=\tfrac{\pi}{2}$, amplitude $\tfrac12$,
        midline $y=-2$, range $[-2.5,-1.5]$.
  \item Period $\pi$; asymptotes where
        $x-\tfrac{\pi}{4}=\tfrac{\pi}{2}+k\pi$, i.e.
        $x=\tfrac{3\pi}{4}+k\pi$.
  \item Amplitude $5$, period $\tfrac{2\pi}{\pi/3}=6$, midline $y=3$,
        shifted $2$ to the right; range $[-2,8]$.
  \item Midline $\tfrac{7+1}{2}=4$, amplitude $\tfrac{7-1}{2}=3$, period
        $4$ so $b=\tfrac{2\pi}{4}=\tfrac{\pi}{2}$:
        $y=3\cos\bigl(\tfrac{\pi x}{2}\bigr)+4$.
  \item $\dfrac{(\cos x-1)(\cos x+1)}{\cos x-1}=\cos x+1$ (for
        $\cos x\neq1$).
  \item $\dfrac{2\sin x\cos x}{\sin 2x}=\dfrac{\sin 2x}{\sin 2x}=1$.
  \item $\dfrac{1}{\cos^{2}x}-\tan^{2}x=1$ by the identity
        $1+\tan^{2}x=1/\cos^{2}x$.
  \item $\sin(x+\pi)=\sin x\cos\pi+\cos x\sin\pi=-\sin x$.
  \item $\cos\bigl(\tfrac{\pi}{2}-x\bigr)
        =\cos\tfrac{\pi}{2}\cos x+\sin\tfrac{\pi}{2}\sin x=\sin x$.
  \item $x=-\tfrac{\pi}{3}+2k\pi$ or
        $x=\pi+\tfrac{\pi}{3}+2k\pi=\tfrac{4\pi}{3}+2k\pi$.
  \item $2x=\tfrac{\pi}{3}+k\pi$, so
        $x=\tfrac{\pi}{6}+\tfrac{k\pi}{2}$.
  \item Factor: $\sin x\,(2\sin x-1)=0$, so $\sin x=0$ giving $x=k\pi$, or
        $\sin x=\tfrac12$ giving $x=\tfrac{\pi}{6}+2k\pi$ and
        $x=\tfrac{5\pi}{6}+2k\pi$.
  \item $2x=\pm\tfrac{\pi}{3}+2k\pi$, so
        $x=\pm\tfrac{\pi}{6}+k\pi$.
  \item No solution: $\sin$ never exceeds $1$, since it is a coordinate of
        a point on the unit circle. An equation can be perfectly
        well-formed and still have an empty solution set.
  \item $\tfrac{\pi}{6}$, \;$\pi$, \;$-\tfrac{\pi}{4}$.
  \item $\sin\tfrac{3\pi}{4}=\tfrac{\sqrt2}{2}$, and
        $\arcsin\tfrac{\sqrt2}{2}=\tfrac{\pi}{4}$. The answer is
        $\tfrac{\pi}{4}$, not $\tfrac{3\pi}{4}$, because $\arcsin$ returns
        values in $[-\tfrac{\pi}{2},\tfrac{\pi}{2}]$ only:
        $\arcsin(\sin x)=x$ holds just on that interval.
  \item Period $\tfrac{2\pi}{\pi/2}=4$ quarters, i.e. one year. Maximum
        $240$ where $\cos\bigl(\tfrac{\pi t}{2}\bigr)=-1$, i.e.
        $\tfrac{\pi t}{2}=\pi+2k\pi$ so $t=2,6,10,\dots$; minimum $160$
        where the cosine is $1$, i.e. $t=0,4,8,\dots$. Demand peaks in the
        third quarter of each year and troughs in the first.
\end{solutions}
